\documentclass[conference]{IEEEtran}
\usepackage{amsmath,amssymb,amsfonts}
\usepackage{algorithmic}
\usepackage{graphicx}
\usepackage{textcomp}
\usepackage{booktabs}
\usepackage{cite}
\usepackage{hyperref}
\usepackage{multirow}
\usepackage{array}
\def\BibTeX{{\rm B\kern-.05em{\sc i\kern-.025em b}\kern-.08em T\kern-.1667em\lower.7ex\hbox{E}\kern-.125emX}}

\begin{document}

\title{Cosine-Normalized Accurate Forgetting for\\Federated Continual Learning}

\author{\IEEEauthorblockN{Sajal Garg}
\IEEEauthorblockA{Department of Computer Science\\Indian Institute of Technology\\Email: sajal@example.edu}}

\maketitle

\begin{abstract}
Federated continual learning (FCL) enables distributed clients to collaboratively learn a sequence of tasks without sharing raw data, yet it faces two compounding challenges: catastrophic forgetting and statistical heterogeneity across clients. The Accurate Forgetting framework (AF-FCL) addresses these issues through conditional normalizing flows that model per-class feature distributions and enable generative replay alongside knowledge distillation. However, the standard linear classification head computes logits as $z_c = \mathbf{W}_c^\top \mathbf{x} + b_c$, coupling both magnitude and direction of weight vectors with the decision. During sequential task updates, this coupling introduces magnitude bias: newly updated class prototypes acquire disproportionately large norms, causing the classifier to favor recent classes irrespective of angular proximity. We propose replacing the linear classifier with a cosine-normalized classification head that computes logits as $z_c = \sigma \cdot \cos(\theta_c)$, where $\theta_c$ is the angle between L2-normalized feature and class prototype vectors, and $\sigma$ is a learnable temperature. This eliminates magnitude bias, produces angularly uniform decision boundaries, and stabilizes normalizing flow training. Experiments on EMNIST-Letters under Non-IID and Shared partitions show the cosine classifier improves class-so-far accuracy from 27.11\% to 62.80\% at the first Non-IID task, achieves 85.25\% final per-task accuracy versus 67.34\% baseline, and reduces flow loss median from 96,449 to 518. On Shared splits, it reaches 87.44\% versus 61.27\% baseline.
\end{abstract}

\begin{IEEEkeywords}
Federated Learning, Continual Learning, Cosine Classifier, Normalizing Flows, Catastrophic Forgetting
\end{IEEEkeywords}

\section{Introduction}

Federated learning (FL) has emerged as a dominant paradigm for privacy-preserving collaborative machine learning, enabling multiple clients to jointly train a shared model without centralizing sensitive data. Each client trains on its local dataset and communicates only model updates to a central server, which aggregates them into a global model.

Continual learning (CL) addresses the problem of learning a sequence of tasks over time without forgetting previously acquired knowledge. The central challenge is catastrophic forgetting, wherein optimization for a new task destructively interferes with representations learned for earlier tasks.

Federated continual learning (FCL) lies at the intersection of these two fields. Distributed clients must learn a sequence of tasks collaboratively, where each client may encounter tasks at different times and with non-identical data distributions. The combination of catastrophic forgetting and statistical heterogeneity creates a particularly challenging optimization landscape.

The Accurate Forgetting framework (AF-FCL) addresses this through a conditional normalizing flow that learns the feature distribution $p(\mathbf{x} \mid y)$ for each class. By sampling from this distribution, synthetic replay features are generated for earlier classes, enabling rehearsal without storing raw data. Combined with knowledge distillation from both the previous local model and the global model, AF-FCL balances plasticity and stability.

However, a fundamental limitation persists: the standard linear classifier computes:
\begin{equation}
z_c = \mathbf{W}_c^\top \mathbf{x} + b_c = \|\mathbf{W}_c\| \cdot \|\mathbf{x}\| \cdot \cos(\theta_c) + b_c
\end{equation}
This entangles the norm $\|\mathbf{W}_c\|$, the feature norm $\|\mathbf{x}\|$, and the angular relationship $\cos(\theta_c)$. During continual learning, weight vectors for recent classes grow in magnitude while those for older classes stagnate, creating \textit{magnitude bias}. In federated settings, heterogeneous client distributions amplify this through uneven update intensities after averaging.

We propose replacing the linear head with a cosine-normalized classifier:
\begin{equation}
z_c = \sigma \cdot \frac{\mathbf{W}_c^\top \mathbf{x}}{\|\mathbf{W}_c\| \cdot \|\mathbf{x}\|} = \sigma \cdot \cos(\theta_c)
\end{equation}
where $\sigma$ is a learnable temperature. By L2-normalizing both vectors, we project representations onto the unit hypersphere, eliminating magnitude from the decision boundary.

Our contributions are: (1) We identify magnitude bias as an underexplored source of FCL instability with geometric analysis. (2) We introduce a CosineLinear module with learnable temperature and optional angular margin into AF-FCL. (3) We demonstrate substantial improvements in accuracy, retention, and flow stability on EMNIST-Letters without modifying replay, distillation, or flow architecture.

\section{Related Work}

\subsection{Federated Learning}
Federated learning was formalized as a distributed optimization problem with FedAvg \cite{mcmahan2017} performing local SGD and weighted averaging. FedProx \cite{li2020} addresses heterogeneity through local regularization. Most FL research assumes fixed task distributions, leaving sequential settings largely unaddressed.

\subsection{Continual Learning}
CL methods fall into three families: regularization-based (EWC \cite{kirkpatrick2017}, Synaptic Intelligence \cite{aljundi2018}), replay-based (iCaRL \cite{rebuffi2017}), and architecture-based approaches. AF-FCL uses normalizing flows as the generative replay mechanism.

\subsection{Cosine Classifiers}
ArcFace \cite{deng2019} demonstrated that normalizing features and prototypes onto the unit hypersphere with additive angular margins produces discriminative representations. LUCIR \cite{hou2019} applied cosine normalization to class-incremental learning in centralized settings, motivating our federated investigation.

\section{Background: The AF-FCL Framework}

\subsection{Problem Formulation}
Consider $K$ clients collaborating through a server over $T$ tasks. Task $t$ introduces classes $\mathcal{C}^t$. Client $k$ receives $\mathcal{D}_k^t = \{(\mathbf{x}_i, y_i)\}$ with $y_i \in \mathcal{C}^t$. The objective is to perform well on all classes $\mathcal{C}^{1:t} = \bigcup_{\tau=1}^{t} \mathcal{C}^\tau$ while accessing only $\mathcal{D}_k^t$.

\subsection{Architecture}
AF-FCL has three components: (1) A \textbf{feature extractor} (S-ConvNet or ResNet-18-CBAM) mapping images to $\mathbf{x}_a \in \mathbb{R}^d$, $d{=}512$. (2) A \textbf{classifier head} with hidden layer $\mathbf{x}_b = \text{LeakyReLU}(\mathbf{W}_2 \mathbf{x}_a)$ and logits $\mathbf{z} = \mathbf{W}_c \mathbf{x}_b$. (3) A \textbf{conditional normalizing flow} modeling $\mathbf{x}_a = f_\phi(\mathbf{u}; y)$, $\mathbf{u} \sim \mathcal{N}(\mathbf{0}, \mathbf{I})$.

\subsection{Training Procedure}
Training alternates between flow and classifier phases. The flow maximizes log-likelihood:
\begin{equation}
\mathcal{L}_{\text{flow}} = -\mathbb{E}\!\left[\log p_\phi(\mathbf{x}_a \mid y)\right] + \lambda_{\text{last}} \mathcal{L}_{\text{flow-last}}
\end{equation}

The classifier loss combines cross-entropy, knowledge distillation, and replay:
\begin{equation}
\mathcal{L}_{\text{cls}} = \mathcal{L}_{\text{CE}} + \lambda_f(\mathcal{L}_{\text{feat-last}} + \mathcal{L}_{\text{feat-global}}) + \lambda_o(\mathcal{L}_{\text{out-last}} + \mathcal{L}_{\text{out-global}}) + \lambda_{\text{flow}} \cdot \mathcal{L}_{\text{replay}}
\end{equation}

The replay loss uses probability-weighted synthetic features:
\begin{equation}
\mathcal{L}_{\text{replay}} = \mathbb{E}_{\tilde{\mathbf{x}}_a, \tilde{y}}\!\left[p_{\text{local}}(\tilde{\mathbf{x}}_a) \cdot \mathcal{L}_{\text{CE}}(\tilde{\mathbf{x}}_a, \tilde{y})\right]
\end{equation}

The log-likelihood is computed via the change-of-variables formula. For the flow $f_\phi$ with base distribution $p_U(\mathbf{u}) = \mathcal{N}(\mathbf{0},\mathbf{I}_d)$:
\begin{equation}
\log p_X(\mathbf{x}) = \log p_U(f_\phi^{-1}(\mathbf{x})) + \log\!\left|\det\frac{\partial f_\phi^{-1}}{\partial \mathbf{x}}\right|
\end{equation}

AF-FCL uses $L{=}4$ affine coupling layers. Each layer partitions input $\mathbf{x} = [\mathbf{x}_{1:d/2}, \mathbf{x}_{d/2+1:d}]$ and applies:
\begin{align}
\mathbf{y}_{1:d/2} &= \mathbf{x}_{1:d/2} \\
\mathbf{y}_{d/2+1:d} &= \mathbf{x}_{d/2+1:d} \odot \exp(s(\mathbf{x}_{1:d/2}; y)) + t(\mathbf{x}_{1:d/2}; y)
\end{align}
where $s,t$ are residual networks conditioned on the class label. The log-determinant simplifies to $\sum_j s_j$, making it $O(d)$.

\section{Proposed Method}

\subsection{The Magnitude Bias Problem}
Expanding the linear logit $z_c = \mathbf{W}_c^\top\mathbf{x} + b_c$ in polar form:
\begin{equation}
z_c = \|\mathbf{W}_c\| \cdot \|\mathbf{x}\| \cdot \cos(\theta_c) + b_c
\end{equation}

Consider two classes $a,b$ with $\theta_a = \theta_b$. If $\|\mathbf{W}_a\| > \|\mathbf{W}_b\|$, then $z_a > z_b$ purely from norm difference. In continual learning, gradient updates increase $\|\mathbf{W}_c\|$ for current-task classes while old-class norms stagnate. The decision boundary between classes $a$ and $b$:
\begin{equation}
(\mathbf{W}_a - \mathbf{W}_b)^\top\mathbf{x} = b_b - b_a
\end{equation}
shifts whenever norms change, even if directions $\hat{\mathbf{W}}_a, \hat{\mathbf{W}}_b$ remain fixed.

\subsection{Cosine-Normalized Classifier}
We define the cosine logit as:
\begin{equation}
z_c = \sigma \cdot \hat{\mathbf{W}}_c^\top \hat{\mathbf{x}} = \sigma \cdot \cos(\theta_c)
\end{equation}
where $\hat{\mathbf{W}}_c = \mathbf{W}_c/\|\mathbf{W}_c\|$ and $\hat{\mathbf{x}} = \mathbf{x}/\|\mathbf{x}\|$. The decision boundary becomes:
\begin{equation}
(\hat{\mathbf{W}}_a - \hat{\mathbf{W}}_b)^\top \hat{\mathbf{x}} = 0
\end{equation}
This is a great circle on $\mathbb{S}^{d-1}$, invariant to any scaling of $\mathbf{W}_a$, $\mathbf{W}_b$, or $\mathbf{x}$.

The softmax probability is:
\begin{equation}
p(y{=}c \mid \mathbf{x}) = \frac{\exp(\sigma\cos\theta_c)}{\sum_{j=1}^{C}\exp(\sigma\cos\theta_j)}
\end{equation}

The temperature $\sigma$ is learnable, initialized to 5.0, and clamped to $[1, 30]$.

\subsection{Angular Margin}
For ground-truth class $y_i$, we apply ArcFace-style margin during training:
\begin{equation}
z_{y_i} = \sigma \cdot \cos(\theta_{y_i} + m)
\end{equation}
Using $\cos(\theta+m) = \cos\theta\cos m - \sin\theta\sin m$, since $\cos m < 1$, this reduces the correct-class logit, forcing tighter angular alignment. The inter-class gap satisfies $\theta_{y_i} + m \leq \theta_c$ for $c \neq y_i$.

\subsection{Gradient Analysis}
For a linear classifier, $\frac{\partial\mathcal{L}}{\partial\mathbf{W}_c} = (p_c - \mathbb{1}[y{=}c])\cdot\mathbf{x}$, scaling with $\|\mathbf{x}\|$. For the cosine classifier:
\begin{equation}
\frac{\partial\mathcal{L}}{\partial\mathbf{W}_c} = \frac{\sigma(p_c - \mathbb{1}[y{=}c])}{\|\mathbf{W}_c\|}\left(\hat{\mathbf{x}} - \cos\theta_c\cdot\hat{\mathbf{W}}_c\right)
\end{equation}
This is inversely proportional to $\|\mathbf{W}_c\|$ (self-regulating) and the direction $(\hat{\mathbf{x}} - \cos\theta_c\hat{\mathbf{W}}_c)$ is purely angular---the tangent component on the hypersphere.

\subsection{Integration with AF-FCL}
Only the final \texttt{nn.Linear} is replaced with \texttt{CosineLinear}. The feature extractor, normalizing flow, replay mechanism, and all distillation losses remain identical. The implementation uses a \texttt{CosineMixin} class that swaps the head post-construction via Python's MRO.

\subsection{EWC Regularization}
After each task $t$, the diagonal Fisher Information Matrix is computed:
\begin{equation}
\hat{F}_{ii}^{(t)} = \frac{1}{N}\sum_{n=1}^{N}\left(\frac{\partial\mathcal{L}_n}{\partial\theta_i}\right)^2
\end{equation}
The EWC penalty is:
\begin{equation}
\mathcal{L}_{\text{EWC}} = \frac{\lambda_{\text{EWC}}}{2}\sum_i F_{ii}^{(t)}(\theta_i - \theta_i^{(t)*})^2
\end{equation}
Fisher matrices across tasks are merged via exponential moving average with $\alpha{=}0.5$.

\subsection{Federated Averaging on the Hypersphere}
FedAvg computes $\mathbf{W}^{\text{global}} = \sum_k \frac{n_k}{n}\mathbf{W}^{(k)}$. With cosine normalization, the forward pass re-normalizes:
\begin{equation}
\hat{\mathbf{W}}_c^{\text{global}} = \frac{\mathbf{W}_c^{\text{global}}}{\|\mathbf{W}_c^{\text{global}}\|}
\end{equation}
Heterogeneous client norms are absorbed, and only directional consensus determines classification.

\section{Experiments}

\subsection{Setup}
We evaluate on EMNIST-Letters (26 classes, 28$\times$28 grayscale) with 10 clients and 6 incremental tasks under Non-IID and Shared splits. The feature extractor is S-ConvNet (channel size 64, $d{=}512$). The flow uses 4 affine coupling layers. Training: 40 global rounds/task, 80 local epochs, Adam ($\text{lr}{=}10^{-3}$), batch size 64. Loss weights: $\lambda_{\text{flow}}{=}0.5$, $\lambda_{\text{KD-last}}{=}0.2$, $\lambda_f{=}0.5$, $\lambda_o{=}0.1$. Cosine settings: $\sigma_{\text{init}}{=}5.0$, $m{=}0.15$, feature calibration enabled, $\lambda_{\text{EWC}}{=}500$.

\subsection{Results: Non-IID Split}

\begin{table}[htbp]
\caption{Non-IID Class-So-Far Accuracy (\%)}
\centering
\begin{tabular}{lccc}
\toprule
Task & Baseline-CPU & Baseline-GPU & \textbf{Cosine} \\
\midrule
t=1 & 27.11 & 26.73 & \textbf{62.80} \\
t=2 & 33.94 & 34.16 & \textbf{49.16} \\
t=3 & 39.28 & 49.11 & \textbf{46.75} \\
t=4 & 41.04 & 40.84 & 33.49 \\
t=5 & 42.46 & 35.01 & \textbf{35.80} \\
t=6 & 36.12 & 35.27 & \textbf{42.73} \\
\bottomrule
\end{tabular}
\end{table}

The cosine classifier achieves 62.80\% at t=1, more than doubling the baseline (27.11\%), and retains advantage at the final task (42.73\% vs 36.12\%).

\begin{table}[htbp]
\caption{Non-IID Final Per-Task Accuracy (\%)}
\centering
\begin{tabular}{lccc}
\toprule
Task & Baseline-CPU & Baseline-GPU & \textbf{Cosine} \\
\midrule
T1 & 27.52 & 24.45 & \textbf{35.48} \\
T2 & 15.30 & 12.53 & \textbf{17.81} \\
T3 & 38.30 & 38.70 & \textbf{40.13} \\
T4 & 29.34 & 29.67 & \textbf{33.69} \\
T5 & 38.90 & 39.48 & \textbf{44.02} \\
T6 & 67.34 & 66.81 & \textbf{85.25} \\
\bottomrule
\end{tabular}
\end{table}

The cosine model outperforms on every task, with T6 showing an 18-point gain (85.25\% vs 67.34\%).

\begin{table}[htbp]
\caption{Non-IID Flow Loss Stability}
\centering
\begin{tabular}{lccc}
\toprule
Config & Median & Min & Max \\
\midrule
Baseline-CPU & 96,449 & 518 & $1.00\times10^{18}$ \\
Baseline-GPU & 20,998 & 521 & $1.53\times10^{18}$ \\
\textbf{Cosine} & \textbf{518} & 470 & $8.38\times10^{11}$ \\
\bottomrule
\end{tabular}
\end{table}

The cosine variant reduces median flow loss by two orders of magnitude (518 vs 96,449) and maximum by seven orders ($10^{11}$ vs $10^{18}$).

\subsection{Results: Shared Split}

\begin{table}[htbp]
\caption{Shared-Split Class-So-Far Accuracy (\%)}
\centering
\begin{tabular}{lcccc}
\toprule
Task & Baseline & Cos-30ep & Cos-50ep & Cos-150ep \\
\midrule
t=1 & 55.95 & 67.27 & 70.43 & \textbf{75.58} \\
t=2 & 54.70 & 71.85 & 74.00 & \textbf{75.67} \\
t=3 & 49.66 & \textbf{75.06} & 73.66 & 74.96 \\
t=4 & 49.50 & 69.65 & \textbf{70.89} & 62.02 \\
t=5 & 52.87 & \textbf{75.69} & 72.36 & 70.76 \\
t=6 & 61.27 & \textbf{87.44} & 86.41 & 84.40 \\
\bottomrule
\end{tabular}
\end{table}

\begin{table}[htbp]
\caption{Shared Final Per-Task Accuracy (\%)}
\centering
\begin{tabular}{lcccc}
\toprule
Task & Baseline & Cos-30ep & Cos-50ep & Cos-150ep \\
\midrule
T1 & 60.55 & 85.62 & 84.73 & \textbf{89.02} \\
T2 & 62.98 & 85.70 & \textbf{86.49} & 85.70 \\
T3 & 54.39 & \textbf{84.98} & 84.65 & 79.56 \\
T4 & 63.02 & \textbf{91.77} & 89.22 & 84.49 \\
T5 & 54.86 & \textbf{86.03} & 81.75 & 75.35 \\
T6 & 71.83 & 90.52 & 91.59 & \textbf{92.24} \\
\bottomrule
\end{tabular}
\end{table}

\begin{table}[htbp]
\caption{Shared Flow Loss Stability}
\centering
\begin{tabular}{lccc}
\toprule
Config & Median & Min & Max \\
\midrule
Baseline & 28,209 & 514 & $1.71\times10^{17}$ \\
Cosine (30ep) & \textbf{470} & 470 & $5.26\times10^{2}$ \\
Cosine (50ep) & \textbf{470} & 470 & $1.42\times10^{5}$ \\
Cosine (150ep) & 642 & 470 & $1.77\times10^{13}$ \\
\bottomrule
\end{tabular}
\end{table}

The cosine variant at 30 epochs achieves 87.44\% final class-so-far accuracy (vs 61.27\% baseline) with near-perfect flow stability (median 470, max 526).

\section{Discussion}

\subsection{Why Cosine Normalization Works in FCL}
In centralized learning, magnitude bias is self-correcting through balanced gradient updates. In continual learning, old classes are absent; in federated learning, different clients update different classes. The cosine classifier eliminates this failure mode entirely by removing magnitude from the decision function.

\subsection{Interaction with Normalizing Flow}
The L2 normalization implicitly encourages stable feature norms, providing the flow with a more stationary input distribution. The cosine variant's median flow loss (470--642) is near the global minimum (470), indicating consistent near-optimal density estimation.

\subsection{Advantages}
The modification is \textbf{lightweight} (adds one parameter $\sigma$, removes bias), \textbf{modular} (no changes to flow/replay/distillation), and \textbf{scalable} (benefits increase with more classes).

\subsection{Limitations}
Performance depends on feature quality---angular proximity must correspond to semantic similarity. The temperature $\sigma$ and margin $m$ require initialization tuning.

\section{Conclusion}

We presented a cosine-normalized classification head for AF-FCL that eliminates magnitude bias by computing logits based purely on angular similarity. On EMNIST-Letters, this improves accuracy (62.80\% vs 27.11\% at first Non-IID task; 87.44\% vs 61.27\% on Shared), reduces forgetting (85.25\% vs 67.34\% final task), and stabilizes flow training (median loss 518 vs 96,449). Future work includes CIFAR-100 evaluation, adaptive margin schedules, and theoretical convergence analysis on the unit hypersphere.

\begin{thebibliography}{14}
\bibitem{mcmahan2017} B.~McMahan \textit{et al.}, ``Communication-efficient learning of deep networks from decentralized data,'' in \textit{Proc. AISTATS}, 2017.
\bibitem{li2020} T.~Li \textit{et al.}, ``Federated optimization in heterogeneous networks,'' in \textit{Proc. MLSys}, 2020.
\bibitem{kirkpatrick2017} J.~Kirkpatrick \textit{et al.}, ``Overcoming catastrophic forgetting in neural networks,'' \textit{PNAS}, vol.~114, no.~13, pp.~3521--3526, 2017.
\bibitem{aljundi2018} R.~Aljundi \textit{et al.}, ``Memory aware synapses: Learning what (not) to forget,'' in \textit{Proc. ECCV}, 2018.
\bibitem{rebuffi2017} S.-A.~Rebuffi \textit{et al.}, ``iCaRL: Incremental classifier and representation learning,'' in \textit{Proc. CVPR}, 2017.
\bibitem{qi2024} Y.~Qi \textit{et al.}, ``Accurate forgetting for federated continual learning,'' \textit{arXiv preprint}, 2024.
\bibitem{deng2019} J.~Deng \textit{et al.}, ``ArcFace: Additive angular margin loss for deep face recognition,'' in \textit{Proc. CVPR}, 2019.
\bibitem{hou2019} S.~Hou \textit{et al.}, ``Learning a unified classifier incrementally via rebalancing,'' in \textit{Proc. CVPR}, 2019.
\bibitem{kingma2015} D.~P.~Kingma and J.~Ba, ``Adam: A method for stochastic optimization,'' in \textit{Proc. ICLR}, 2015.
\bibitem{dinh2017} L.~Dinh \textit{et al.}, ``Density estimation using real-valued non-volume preserving transformations,'' in \textit{Proc. ICLR}, 2017.
\bibitem{papamakarios2021} G.~Papamakarios \textit{et al.}, ``Normalizing flows for probabilistic modeling and inference,'' \textit{JMLR}, vol.~22, no.~57, pp.~1--64, 2021.
\bibitem{chaudhry2018} A.~Chaudhry \textit{et al.}, ``Riemannian walk for incremental learning,'' in \textit{Proc. ECCV}, 2018.
\bibitem{li2018} Z.~Li and D.~Hoiem, ``Learning without forgetting,'' \textit{IEEE TPAMI}, vol.~40, no.~12, pp.~2935--2947, 2018.
\bibitem{zenke2017} F.~Zenke \textit{et al.}, ``Continual learning through synaptic intelligence,'' in \textit{Proc. ICML}, 2017.
\end{thebibliography}

\end{document}
